O objetivo deste trabalho é estudar a convexidade $P3$ orientada e alguns parâmetros relacionados a mesma. Inicialmente é apresentada uma revisão dos conceitos básicos sobre convexidade em grafos e conceitos básicos sobre teoria dos grafos no geral. Posteriormente, é apresentado um estudo sobre o problema do \emph{Espectro} do número de intervalo, onde verifica-se a continuidade do mesmo, ou seja, todos os valores inteiros entre os limitantes fazem parte do espectro e em seguida, é calculado alguns limitantes para algumas classes de grafos. Além de resultados estruturais, são apresentados algoritmos eficientes para obter os parâmetros de número de intervalo e número de envoltória em árvores orientadas. Por fim, este trabalho apresenta também, resultados de complexidade para a classe de grafos planares orientados, onde temos que os parâmetros de número de intervalo e número de envoltória são ambos \npcs.

% Separe as palavras-chave por ponto
\palavraschave{Convexidade em Grafos. Grafos orientados. Número de intervalo. Número de envoltória. Complexidade computacional.}